% Laget automatisk av elevark/lag_elevark.py fra sim/linearisering.html. Endringer her blir overskrevet.
\documentclass[11pt]{article}
\usepackage{elevark}
\begin{document}
\elevarkhode{Fysikk 2 · Tenke og gjøre fysikk}{Linearisering og modeller}{https://fysikklab.no/sim/linearisering.html}{Simuleringen}

En fysisk modell sier hvordan to størrelser henger sammen, for eksempel $T = 2π\sqrt{L/g}$. Om modellen stemmer, kan vi teste ved å velge akser slik at modellen gir en rett linje. Ligger målepunktene på linja, styrker det modellen, og stigningstallet gir oss konstanten.

\section*{Oppgaver}

\begin{oppgaveliste}
\oppgave \fase{gjett}{Gjett}Se på pendeldataene med $T$ mot $L$. Er grafen en rett linje? Hvilke akser må du velge for at modellen $T = 2π\sqrt{L/g}$ skal gi en rett linje gjennom origo? Test.
\svar
\oppgave \fase{test}{Regn}Vis at $T^{2} = (4π^{2}/g) \cdot L$. Bruk stigningstallet til å bestemme $g$.
\svar
\oppgave \fase{test}{Undersøk}Fritt fall: Hvilke akser gir en rett linje? Bestem $g$. Hvorfor er ikke $R^{2}$ alene nok til å avgjøre hvilken modell som er riktig?
\svar
\oppgave \fase{test}{Undersøk}Magnetfelt rundt en leder med strømmen 5,0~A: Finn transformasjonen som gir en rett linje gjennom origo, og bestem $μ_{0}$. Sammenlign med $4π \cdot 10^{-7}\,\mathrm{T\,m/A}$.
\svar
\oppgave \fase{forklar}{Forklar}Muffinsformer: Prøv både $v_{\mathrm{t}}$ mot $n$ og $v_{\mathrm{t}}^{2}$ mot $n$. Hvilken modell for luftmotstanden støtter dataene, $L \propto v$ eller $L \propto v^{2}$? Begrunn svaret.
\svar
\oppgave \fase{forklar}{Forklar}Hva betyr det at en modell blir «styrket» og ikke «bevist» av målingene? Hva skal til for å forkaste en modell?
\svar
\oppgave \fase{utfordring}{Utfordring}Lim inn egne data fra et forsøk, for eksempel pendel eller muffinsformer. Finn transformasjonen som gir best rett linje, og bestem konstanten.
\svar
\end{oppgaveliste}

\end{document}
